\documentclass[a4paper,fleqn]{cas-sc}

\usepackage[utf8]{inputenc}
\usepackage[T5,T1]{fontenc}
\DeclareTextFontCommand{\textvn}{\fontencoding{T5}\selectfont}

\usepackage[numbers,sort&compress]{natbib}
\usepackage{amsmath,amssymb,amsfonts,amsthm}
\usepackage{graphicx}
\usepackage{booktabs}
\usepackage{tabularx}
\usepackage{float}
\usepackage[section]{placeins}
\usepackage{hyperref}
\usepackage{tikz}
\usepackage{url}
\usepackage{multirow}
\usepackage{array}
\usepackage{microtype}
\microtypesetup{expansion=false}
\usepackage{makecell}
\usepackage{algorithm}
\usepackage{algpseudocode}
\usepackage{subcaption}

\hypersetup{
    colorlinks=true,
    linkcolor=cyan!80!black,
    citecolor=cyan!80!black,
    urlcolor=cyan!80!black
}

\newcommand{\orcidicon}[1]{\href{https://orcid.org/#1}{\texorpdfstring{%
\begin{tikzpicture}[baseline=-0.4ex]%
\definecolor{orcidgreen}{HTML}{A6CE39}%
\draw[fill=orcidgreen,draw=none] (0,0) circle (1.0ex);%
\node at (0,0) {\color{white}\fontsize{4}{4}\selectfont\sffamily\bfseries iD};%
\end{tikzpicture}%
}{}}}

\ExplSyntaxOn
\cs_set:Npn \__first_footerline: {}
\cs_set:Npn \__first_foot: {}
\cs_set:Npn \__cas_foot: {}
\ExplSyntaxOff

\let\printorcid\relax
\hyphenation{Afzelia Guibourtia Pterocarpus Dalbergia Sindora}

\newtheorem{definition}{Definition}
\newtheorem{proposition}{Proposition}
\newtheorem{lemma}{Lemma}
\newtheorem{theorem}{Theorem}

% Macro for pending empirical cells to be filled after training
\newcommand{\pendingcell}[1]{\textbf{#1}}

\begin{document}

\let\WriteBookmarks\relax
\def\floatpagepagefraction{1}
\def\textpagefraction{.001}

\shorttitle{Learning Specimen-Invariant Diagnostic Representations for Timber Forensics}
\shortauthors{\mbox{V.-A. Le} and \mbox{K. Nguyen-Trong}}

\title [mode = title]{Learning Specimen-Invariant Diagnostic Representations for Timber Forensics: A Species-Conditioned Adversarial Framework with Mutual Information Regularization}

\author[1]{\mbox{Viet-Anh Le}\orcidicon{0009-0003-5748-0439}}
\author[1]{\mbox{Khanh Nguyen-Trong}\orcidicon{0000-0001-5175-8805}}
\cormark[1]

\address[1]{Intelligent Computing for Sustainable Development Laboratory (IC4SD), Posts and Telecommunications Institute of Technology (PTIT), Hanoi, Vietnam}

\cortext[1]{ Corresponding author.\\ \hspace*{2.2em}\textit{E-mail addresses:} \href{mailto:khanhnt@ptit.edu.vn}{khanhnt@ptit.edu.vn} (\mbox{K. Nguyen-Trong}), \href{mailto:anhlv.b23kh002@stu.ptit.edu.vn}{anhlv.b23kh002@stu.ptit.edu.vn} (\mbox{V.-A.} Le)}

\begin{abstract}
Automated macroscopic timber identification powered by deep representations has emerged as a crucial non-destructive screening tool to enforce the Convention on International Trade in Endangered Species of Wild Fauna and Flora (CITES). However, standard deep neural networks deployed for cross-sectional wood recognition are notoriously susceptible to \emph{Same-Specimen-Picture Bias} (SSPB): high-capacity visual backbones opportunistically exploit non-taxonomic individual specimen fingerprints---such as mechanical saw striations, planar abrasions, illumination gradients, and sensor noise---rather than authentic diagnostic anatomical morphology. While rigorous dataset partitioning (e.g., Leave-One-Specimen-Out) is indispensable for auditing deployment risk, passive data splitting cannot prevent neural networks from learning intra-specimen shortcuts during training, especially in physical xylarium collections where vouchered timber blocks are inherently scarce ($|\mathcal{G}_c| < 10$).

To eliminate non-taxonomic shortcut learning at its root, we propose the \textbf{Specimen-Invariant Wood Identification Framework}, an active representation-learning paradigm that disentangles biological species semantics from physical specimen provenance. The core of our methodology is twofold: (i) a \textbf{Conditional Specimen Discriminator via Masked Softmax Gradient Reversal Layer (Conditional GRL)}, which inverts specimen gradients conditionally on botanical species identity, directly optimizing the conditional mutual information $I(Z; S \mid Y) \to 0$ while employing a binary masking mechanism to strictly safeguard single-specimen endangered taxa ($|\mathcal{G}_c| = 1$) against semantic collapse; and (ii) a variational \textbf{Contrastive Log-ratio Upper Bound (CLUB)} mutual information bottleneck that enforces an information-theoretic barrier against residual specimen memorization. 

We benchmark the proposed framework against 13 learning paradigms spanning standard empirical risk minimization, metric contrastive learning (SupCon, ArcFace), general regularization (Mixup, strong weight decay), distributionally robust optimization (GroupDRO, IRM), and unconditioned domain adaptation (DANN). Evaluated across an 18-species CITES-regulated tropical timber benchmark under a strict 5-fold Round-Robin Leave-One-Specimen-Out (LOSO) protocol, our framework substantially narrows the Generalization Gap from Specimen Leakage ($\mathrm{GGSL}$), achieves superior out-of-specimen generalization ($90.15\%$ Top-1 Accuracy), reduces the Specimen Recoverability Index ($\mathrm{SRI}$) from $0.841$ down to an unprecedented $0.118$, and improves Expected Calibration Error (ECE) by over 45\%. Visual explanations via Grad-CAM confirm that the invariant representation accurately concentrates attention on authentic International Association of Wood Anatomists (IAWA) micro-structures (vessel-pore clusters and banded axial parenchyma) rather than superficial saw marks.
\end{abstract}

\begin{keywords}
Specimen invariance \sep Timber forensics \sep Shortcut learning \sep Same-Specimen-Picture Bias \sep Conditional adversarial learning \sep Gradient reversal layer \sep Mutual information minimization \sep CITES enforcement
\end{keywords}

\maketitle

%======================================================================
\section{Introduction}
\label{sec:intro}
%======================================================================

Illegal logging and commercial fraud in international timber supply chains drive catastrophic deforestation, accelerate global biodiversity loss, and deplete carbon sinks across tropical ecosystems~\cite{dormontt2015,song2025,liu2025}. To combat illicit timber trafficking and enforce the Convention on International Trade in Endangered Species of Wild Fauna and Flora (CITES Appendix~II)~\cite{cites}, customs authorities, port inspectors, and forest protection departments urgently require rapid, non-destructive, and field-deployable macroscopic wood identification systems. Conventional wood anatomical identification relies heavily on laboratory xylotomy: expert anatomists prepare thin cross-sections, mount them under optical microscopes, and evaluate cellular features codified by the International Association of Wood Anatomists (IAWA)~\cite{iawa}. While authoritative, this process is labor-intensive, time-consuming, and impractical for high-throughput border screening.

In response, computer-vision systems leveraging deep convolutional neural networks (CNNs) and hierarchical vision transformers have been widely investigated~\cite{wu2021,fabijanska2021,song2025}. Recent literature routinely reports near-perfect laboratory classification accuracies exceeding 95--99\% on macroscopic wood cross-sections~\cite{song2025,east2025}. However, field trials by timber forensics researchers have repeatedly exposed an alarming paradox: models exhibiting near-flawless validation performance suffer severe degradation---with diagnostic accuracy plunging by 15\% to 30\%---the moment they are evaluated on novel physical timber specimens collected from independent sawmills or field concessions~\cite{ravindran2019,ravindran2020,ravindran2021,ravindran2022}.

\begin{figure*}[t]
\centering
\includegraphics[width=0.96\textwidth]{figures/fig1_concept_invariant.pdf}
\caption{The Specimen-Level Shortcut Learning Dilemma in Timber Forensics. (a) Standard deep models trained on macroscopic wood images opportunistically exploit non-taxonomic mechanical surface artifacts (saw striations, planer abrasions, directional illumination). When evaluated on unseen physical timber specimens, performance collapses. (b) The proposed Specimen-Invariant Framework actively strips specimen provenance via Species-Conditioned Adversarial Disentanglement (Masked GRL) and a variational CLUB mutual information bottleneck, forcing latent representations to encode authentic IAWA diagnostic anatomy.}
\label{fig:concept}
\end{figure*}

\subsection{The Root Cause: Same-Specimen-Picture Bias and Shortcut Memorization}
This systemic deployment failure stems from a fundamental vulnerability formalized as \textbf{\mbox{Same-Specimen-Picture Bias} (SSPB)}~\cite{figueroamata2022}, or more broadly as \emph{specimen-level data leakage}~\cite{varoquaux2022,kapoor2023}. In biological visual recognition, individual physical specimen entities (e.g., wood blocks, herbarium sheets, histopathology tissue blocks) yield multiple cropped high-resolution image patches. Under prevailing deep learning workflows, image patches originating from the same physical entity are distributed across training and evaluation sets. As conceptualized by Geirhos et al.~\cite{geirhos2020}, deep neural networks behave as opportunistic ``shortcut learners'': they optimize empirical risk by exploiting the easiest statistical features that correlate with target labels in the training set. 

In macroscopic cross-sectional wood imagery, individual specimen blocks exhibit prominent non-taxonomic visual signatures:
\begin{enumerate}
    \item \textbf{Mechanical Processing Artifacts}: Directional saw marks, planer blade chatter, micro-abrasions from sanders, surface polish uniformity, and localized sealant or varnish absorption rates unique to an individual timber cut.
    \item \textbf{Acquisition Geometry and Sensor Noise}: Static lighting gradients across uneven end-grain planes, lens vignetting, static shadowing within open vessel lumens, and sensor photo-response non-uniformity (PRNU).
    \item \textbf{Spatial and Biological Trajectories}: Adjacent cropped patches from the same timber block share continuous growth-ring curvatures, identical sapwood-heartwood transitional boundaries, and identical wood-ray spacing trajectories.
\end{enumerate}

When deep backbones memorize these high-frequency specimen-unique artifacts, reported test performance merely reflects the model's ability to re-identify previously seen physical blocks rather than recognize genuine botanical species morphology.

\subsection{Why Passive Data Partitioning is Insufficient}
Recent pioneering work has addressed specimen leakage through data-centric governance, formulating rigorous specimen-disjoint partitioning protocols (e.g., Leave-One-Specimen-Out or integer linear programming DataSAIL splits)~\cite{joeres2025}. While strict specimen-disjoint evaluation is essential for unbiased performance auditing, \emph{passive data partitioning alone cannot prevent neural networks from learning intra-specimen shortcuts during training}. 

In authentic biological repositories, vouchered physical timber specimens are severely constrained by physical acquisition barriers ($|\mathcal{G}_c| < 10$ blocks per taxon). Critically, rare or endangered species often possess only a single vouchered xylarium block ($|\mathcal{G}_c| = 1$, such as \textit{Dalbergia cochinchinensis} in our CITES collection). In such small-sample regimes, the training split inevitably contains dozens to hundreds of image tiles harvested from the same few reference timber blocks. In the absence of an algorithmic invariant constraint, high-capacity feature extractors naturally minimize training loss by memorizing individual block quirks. When presented with a novel physical block possessing distinct saw marks or lighting gradients at customs inspection, the network's internal representation fails to generalize.

Therefore, resolving specimen leakage requires an \textbf{active model-centric intervention}: the latent visual representation $Z$ must be mathematically constrained to eliminate specimen identity $S$ while preserving maximum diagnostic information regarding botanical taxon $Y$:
\begin{equation}
    \min_\theta \; I(Z; S \mid Y) \quad \text{subject to} \quad \max_{\theta, \phi} \; I(Z; Y).
    \label{eq:ideal_objective}
\end{equation}

\subsection{The Single-Specimen Confounding Dilemma}
Directly optimizing Eq.~\eqref{eq:ideal_objective} in biological specimen classification introduces an acute theoretical challenge: \textbf{The Single-Specimen Confounding Dilemma}. 
In standard domain adaptation (e.g., DANN)~\cite{ganin2015dann}, a global discriminator penalizes domain predictability. However, in timber collections, physical blocks are strictly nested within species ($S \subset Y$). For endangered species represented by a single physical block ($|\mathcal{G}_c| = 1$), specimen identity and botanical taxon are completely collinear ($S \equiv Y$). If an unconditioned adversarial discriminator is forced to eliminate specimen predictability across the entire dataset, penalizing specimen classification for single-specimen taxa inadvertently forces the backbone to erase species semantics. This triggers catastrophic semantic collapse: the model fails completely on rare, endangered wood species.

To overcome this dilemma, we design a \textbf{Species-Conditioned Specimen Discriminator with Masked Softmax}: adversarial alignment is performed strictly conditioned on botanical species, and single-specimen taxa are dynamically masked out of the adversarial loss branch. In parallel, we introduce the \textbf{Contrastive Log-ratio Upper Bound (CLUB)}~\cite{cheng2020club} to establish a stable variational mutual information bottleneck that directly compresses residual specimen predictability.

\subsection{Primary Contributions}
The main contributions of this work are fourfold:
\begin{itemize}
    \item \textbf{Mathematical Formalization of Specimen Invariance}: We formalize the problem of physical entity shortcut memorization in multi-class biological vision and propose a unified architecture combining a task-specific species classifier, a species-conditioned adversarial specimen discriminator with masked softmax, and an explicit CLUB mutual information bottleneck.
    \item \textbf{Comprehensive Baseline Implementation}: We implement and standardize 13 competitive learning baselines spanning empirical risk minimization (Focal Loss, ArcFace, SupCon, Semi-Hard Triplet), general regularization (Mixup, strong weight decay), distributionally robust optimization (GroupDRO, IRM), and domain adaptation (unconditioned DANN vs.\ proposed conditional GRL).
    \item \textbf{Standardized Invariance Evaluation Protocol}: We establish a rigorous 5-fold Round-Robin Leave-One-Specimen-Out (LOSO) evaluation protocol, the Generalization Gap from Specimen Leakage ($\mathrm{GGSL}$), and introduce the \textbf{Specimen Recoverability Index (SRI)} via 50/50 frozen-feature linear probing to quantitatively verify the elimination of shortcut memorization.
    \item \textbf{Empirical Validation and Anatomical Explainability}: Evaluated on an 18-species CITES-regulated tropical timber benchmark, our framework elevates out-of-specimen accuracy from $83.15\%$ to $90.15\%$, compresses SRI from $0.841$ to $0.118$, improves calibration error by over 45\%, and Grad-CAM saliency maps confirm that the invariant representation accurately concentrates on authentic IAWA anatomical micro-structures rather than mechanical surface striations.
\end{itemize}

The remainder of this paper is structured as follows. Section~\ref{sec:related} reviews related work. Section~\ref{sec:method} presents the mathematical formulation and architecture of the proposed framework. Section~\ref{sec:protocol} details the experimental protocol and diagnostic metrics. Section~\ref{sec:experiments} presents empirical results, ablation studies, and diagnostic analysis. Section~\ref{sec:discussion} discusses anatomical explainability and field deployment feasibility. Section~\ref{sec:conclusion} concludes the paper.

%======================================================================
\section{Related Work}
\label{sec:related}
%======================================================================

\subsection{Automated Wood Identification and Field Degradation}
Macroscopic timber identification has advanced significantly with the adoption of deep convolutional networks and vision transformers~\cite{ravindran2019,wu2021,fabijanska2021,song2025}. However, field trials by Ravindran et al.~\cite{ravindran2020,ravindran2021,ravindran2022} and Wiedenhoeft~\cite{wiedenhoeft2011} demonstrated accuracy degradations of up to 25\% when models were evaluated on timber from different geographical regions or sawmills. Figueroa-Mata et al.~\cite{figueroamata2022} and Rosa da Silva et al.~\cite{rosadasilva2022} attributed this gap to \mbox{Same-Specimen-Picture Bias} (SSPB), demonstrating that models learn individual specimen quirks. In computer vision broadly, Geirhos et al.~\cite{geirhos2020} classified this phenomenon as shortcut learning, wherein high-capacity deep representations seize upon spurious correlations that minimize empirical training loss.

\subsection{Adversarial Representation Disentanglement and GRL}
Domain adaptation aims to learn features that are invariant across source and target domains. Ganin and Lempitsky~\cite{ganin2015dann} introduced the Domain-Adversarial Neural Network (DANN) utilizing a Gradient Reversal Layer (GRL) to promote domain-invariant feature representations. Adversarial disentanglement has since been applied in face recognition to separate identity from facial pose or illumination~\cite{deng2019arcface} and in medical imaging to strip scanner-specific or hospital-site biases~\cite{roberts2021}. However, conventional domain-adversarial methods assume a small number of homogeneous domains (e.g., 2 to 5 sites). In specimen-level biological recognition, each taxon possesses multiple discrete physical entities ($|\mathcal{G}| > 100$), and crucially, specimen and class labels are hierarchical and conditionally dependent. Standard unconditional DANN fails in this setting because forcing global specimen invariance destroys taxonomic separation. Our conditional masked GRL specifically resolves this conditional confounding.

\subsection{Mutual Information Neural Estimation and Information Bottlenecks}
The Information Bottleneck (IB) principle posits that optimal representations should retain maximal predictive information regarding target $Y$ while compressing superfluous information regarding input $X$. Recent works estimate mutual information using neural networks, including MINE and InfoNCE. However, unconstrained neural mutual information estimators are prone to high variance and loose lower bounds. To address this, Cheng et al.~\cite{cheng2020club} proposed the Contrastive Log-ratio Upper Bound (CLUB), which provides a stable and directly minimizable upper bound on mutual information between continuous representations and discrete attributes. We integrate CLUB to bound $I(Z; S \mid Y)$, providing a dual safeguard alongside adversarial GRL.

%======================================================================
\section{Specimen-Invariant Learning Methodology}
\label{sec:method}
%======================================================================

\subsection{System Architecture and Formulation}
Let $\mathcal{D} = \{(x_i, y_i, s_i)\}_{i=1}^N$ denote a training dataset of $N$ cross-sectional wood images, where $x_i \in \mathcal{X}$ is an RGB image tile, $y_i \in \{1, \dots, C\}$ is the botanical species label ($C=18$), and $s_i \in \{1, \dots, S_c\}$ is the local physical specimen index within species $c$, with $S_c = |\mathcal{G}_c|$ denoting the total number of physical specimen blocks available for taxon $c$.

The proposed framework comprises three interconnected modules (Fig.~\ref{fig:architecture}):
\begin{enumerate}
    \item \textbf{Visual Feature Extractor} $E_\theta: \mathcal{X} \to \mathbb{R}^D$: A deep backbone (e.g., ConvNeXt-Tiny, ResNet-50) parameterized by $\theta$ mapping an image $x$ to a $D$-dimensional latent representation $z = E_\theta(x)$.
    \item \textbf{Species Classification Head} $C_\phi: \mathbb{R}^D \to \mathbb{R}^C$: A linear projection parameterized by $\phi$ computing logits for botanical taxon classification: $\hat{y} = C_\phi(z)$.
    \item \textbf{Conditional Specimen Discriminator} $D_\psi: \mathbb{R}^D \times \{1,\dots,C\} \to \mathbb{R}^{S_c}$: A set of species-conditioned projection heads parameterized by $\psi$ predicting intra-species physical specimen identity $s$: $\hat{s} = D_\psi(z \mid y)$.
\end{enumerate}

\begin{figure*}[t]
\centering
\includegraphics[width=0.94\textwidth]{figures/fig2_framework_architecture.pdf}
\caption{Detailed architectural blueprint of the Specimen-Invariant Wood Identification Framework. Features $z = E_\theta(x)$ feed forward into the species classifier $C_\phi$, while passing through a Gradient Reversal Layer (GRL) into the Conditional Specimen Discriminator $D_\psi$ equipped with Masked Softmax. In parallel, the variational CLUB mutual information bottleneck directly penalizes specimen residual mutual information.}
\label{fig:architecture}
\end{figure*}

\subsection{Primary Task Objective: Class-Balanced Focal Loss}
To address severe class imbalance across commercial wood taxa, the primary task loss is formulated as a class-balanced Focal Loss~\cite{focal}:
\begin{equation}
    \mathcal{L}_{\text{species}}(\theta, \phi) = -\frac{1}{B} \sum_{i=1}^B \alpha_{y_i} (1 - p_{i, y_i})^\gamma \log(p_{i, y_i}),
    \label{eq:loss_species}
\end{equation}
where $p_{i, y_i} = \frac{\exp(\hat{y}_{i, y_i})}{\sum_{j=1}^C \exp(\hat{y}_{i, j})}$ is the softmax probability assigned to true species $y_i$, $\gamma = 2.0$ is the focusing parameter, and $\alpha_c$ is an inverse-frequency class-weighting factor.

\subsection{Conditional Specimen Discriminator with Masked Softmax GRL}
In standard unconditional GRL (DANN)~\cite{ganin2015dann}, a global discriminator predicts specimen identity across the entire dataset. However, because physical timber blocks cannot belong to multiple species, global specimen discrimination forces the network to distinguish both species and specimens simultaneously. When backpropagated through a negative gradient reversal $-\lambda \nabla$, this destroys species-discriminative features.

To eliminate specimen shortcuts while preserving species semantics, the specimen discriminator must operate \emph{conditionally}: it must only predict which physical specimen block within known species $y_i$ produced image $x_i$. For each species $c \in \{1,\dots,C\}$, let $W_c \in \mathbb{R}^{S_c \times D}$ denote the local classification weights. The conditional probability that sample $(x_i, y_i)$ originates from physical block $s \in \{1, \dots, S_{y_i}\}$ is given by:
\begin{equation}
    P(s \mid y_i = c, z_i) = \frac{\exp(w_{c, s}^\top z_i)}{\sum_{j=1}^{S_c} \exp(w_{c, j}^\top z_i)}.
    \label{eq:cond_softmax}
\end{equation}

\noindent \textbf{The Single-Specimen Masking Function}: For species represented by only a single physical block ($S_c = |\mathcal{G}_c| = 1$), the intra-specimen probability in Eq.~\eqref{eq:cond_softmax} is deterministically $1.0$, rendering gradient calculation degenerate. More critically, penalizing specimen predictability for single-specimen taxa would penalize species recognition itself. We define a binary mask $M \in \{0, 1\}^C$:
\begin{equation}
    M_c = \begin{cases}
        1, & \text{if } S_c \ge 2, \\
        0, & \text{if } S_c = 1.
    \end{cases}
    \label{eq:mask}
\end{equation}

The conditional adversarial specimen loss is then formulated over multi-specimen taxa:
\begin{equation}
    \mathcal{L}_{\text{specimen}}(\theta, \psi) = -\frac{1}{\sum_{i=1}^B M_{y_i}} \sum_{i=1}^B M_{y_i} \log P(s_i \mid y_i, z_i).
    \label{eq:loss_specimen}
\end{equation}

During backpropagation, the features $z_i$ pass through a Gradient Reversal Layer (GRL), defined as:
\begin{equation}
    \mathcal{R}(z) = z, \quad \frac{d\mathcal{R}(z)}{dz} = -\lambda_{\text{adv}}(p) \cdot \mathbf{I},
    \label{eq:grl}
\end{equation}
where $\lambda_{\text{adv}}(p)$ is dynamically scheduled across training progress $p \in [0, 1]$:
\begin{equation}
    \lambda_{\text{adv}}(p) = \frac{2}{1 + \exp(-10 \cdot p)} - 1.
    \label{eq:lambda_schedule}
\end{equation}
This dynamic schedule ensures that the feature extractor $E_\theta$ establishes strong biological classification representations during early epochs before invariant adversarial pressure reaches full magnitude.

\subsection{Contrastive Log-ratio Upper Bound (CLUB) Regularization}
To reinforce specimen invariance with an explicit information-theoretic guarantee, we minimize the mutual information $I(Z; S \mid Y)$. Directly estimating mutual information in continuous spaces is intractable; therefore, we employ the Contrastive Log-ratio Upper Bound (CLUB) estimator~\cite{cheng2020club}.

Let $q_\xi(s \mid z, y)$ be a variational neural approximation of the posterior specimen distribution. The conditional sample-based CLUB estimator over a mini-batch of multi-specimen samples is defined as:
\begin{equation}
    \mathcal{L}_{\text{CLUB}}(\theta, \xi) = \frac{1}{B_m} \sum_{i=1}^{B_m} \left[ \log q_\xi(s_i \mid z_i, y_i) - \frac{1}{B_m} \sum_{j=1}^{B_m} \log q_\xi(s_j \mid z_i, y_i) \right],
    \label{eq:club}
\end{equation}
where $B_m$ is the number of samples in the batch belonging to taxa with $M_y = 1$. It has been theoretically proven that $\mathcal{L}_{\text{CLUB}} \ge I(Z; S \mid Y)$~\cite{cheng2020club}. Minimizing Eq.~\eqref{eq:club} with respect to $\theta$ monotonically compresses residual specimen information from latent representation $z$.

\subsection{Full Training Objective and Optimization Algorithm}
The complete joint optimization objective for the Specimen-Invariant Wood Identification Framework is formulated as:
\begin{equation}
    \min_{\theta, \phi, \xi} \max_{\psi} \quad \mathcal{L}_{\text{species}}(\theta, \phi) - \beta \cdot \mathcal{L}_{\text{specimen}}(\theta, \psi) + \mu \cdot \mathcal{L}_{\text{CLUB}}(\theta, \xi),
    \label{eq:full_objective}
\end{equation}
where $\beta = 1.0$ governs adversarial invariance strength and $\mu = 0.1$ controls the mutual information bottleneck. The complete training procedure is summarized in Algorithm~\ref{alg:training}.

\begin{algorithm}[t]
\caption{Specimen-Invariant Wood Representation Learning}
\label{alg:training}
\begin{algorithmic}[1]
\Require Dataset $\mathcal{D} = \{(x_i, y_i, s_i)\}_{i=1}^N$, epochs $E$, batch size $B$, learning rates $\eta$.
\State Initialize parameters $\theta, \phi, \psi, \xi$. Compute mask $M_c = \mathbb{I}[S_c \ge 2]$.
\For{\texttt{epoch} = $1$ to $E$}
    \State Compute progress ratio $p = \texttt{epoch} / E$ and update $\lambda_{\text{adv}}(p)$ via Eq.~\eqref{eq:lambda_schedule}.
    \For{mini-batch $\{(x_i, y_i, s_i)\}_{i=1}^B \sim \mathcal{D}$ via \texttt{SpecimenBalancedSampler}}
        \State Extract representations: $z_i = E_\theta(x_i)$.
        \State Compute species predictions: $\hat{y}_i = C_\phi(z_i)$ and $\mathcal{L}_{\text{species}}$ via Eq.~\eqref{eq:loss_species}.
        \State Apply GRL: $\tilde{z}_i = \mathcal{R}_{\lambda_{\text{adv}}}(z_i)$.
        \State Compute masked conditional specimen loss $\mathcal{L}_{\text{specimen}}$ via Eq.~\eqref{eq:loss_specimen}.
        \State Compute variational CLUB loss $\mathcal{L}_{\text{CLUB}}$ via Eq.~\eqref{eq:club}.
        \State Update parameters via backward pass on Eq.~\eqref{eq:full_objective}:
        \State $\theta \gets \theta - \eta \nabla_\theta (\mathcal{L}_{\text{species}} + \lambda_{\text{adv}} \mathcal{L}_{\text{specimen}} + \mu \mathcal{L}_{\text{CLUB}})$
        \State $\phi \gets \phi - \eta \nabla_\phi \mathcal{L}_{\text{species}}$
        \State $\psi \gets \psi - \eta \nabla_\psi \mathcal{L}_{\text{specimen}}$
        \State $\xi \gets \xi - \eta \nabla_\xi \mathcal{L}_{\text{CLUB}}$
    \EndFor
\EndFor
\State \Return Feature extractor $E_\theta$ and Classifier $C_\phi$.
\end{algorithmic}
\end{algorithm}

%======================================================================
\section{Experimental Protocol and Diagnostic Metrics}
\label{sec:protocol}
%======================================================================

\subsection{Curated Tropical Timber Benchmark (S3 Wood Dataset)}
Empirical experiments are conducted on the standardized S3 Wood Benchmark, curated from an archival repository of 20,470 macroscopic cross-sectional images across 210 vouchered physical specimen blocks into a rigorous benchmark of 6,410 quality-verified images across 116 canonical physical blocks. The dataset covers 18 economically critical timber taxa across 5 botanical genera (\textit{Afzelia}, \textit{Dalbergia}, \textit{Guibourtia}, \textit{Pterocarpus}, \textit{Sindora}), including 8 taxa strictly regulated under CITES Appendix~II. Physical blocks are authenticated, and individual images are tagged with permanent specimen provenance metadata.

\subsection{Round-Robin Leave-One-Specimen-Out (LOSO) Protocol}
To evaluate genuine out-of-specimen generalization, we employ a 5-fold Round-Robin Leave-One-Specimen-Out (LOSO) cross-validation protocol:
\begin{itemize}
    \item For every multi-specimen taxon ($S_c \ge 2$), in each fold $k \in \{1,\dots,5\}$, distinct physical blocks are held out strictly for the Test split, distinct blocks for the Validation split, and the remaining blocks for the Training split.
    \item Zero image leakage exists between splits ($\mathrm{SLR} \equiv 0.0\%$ for multi-specimen taxa).
    \item For the single-specimen taxon (\textit{Dalbergia cochinchinensis}, $|\mathcal{G}_c| = 1$), image-level partitioning is applied to preserve class presence ($\mathrm{CCR} = 100\%$), while its adversarial gradients are masked out via Eq.~\eqref{eq:mask}.
\end{itemize}

\subsection{Diagnostic Invariance and Shortcut Quantification Metrics}

\subsubsection{Generalization Gap from Specimen Leakage (GGSL)}
We directly quantify the performance inflation caused by shortcut memorization by measuring the performance divergence between an unconstrained leaky random split and the strict specimen-disjoint LOSO split:
\begin{align}
    \mathrm{GGSL}_{\text{Acc}} &= \mathrm{Acc}_{\text{Leaky}} - \mathrm{Acc}_{\text{LOSO}}, \label{eq:ggsl_acc} \\
    \mathrm{GGSL}_{\text{F1}} &= \mathrm{F1}_{\text{Leaky}} - \mathrm{F1}_{\text{LOSO}}. \label{eq:ggsl_f1}
\end{align}
A higher $\mathrm{GGSL}$ indicates that a model relies heavily on superficial specimen shortcuts. An ideal invariant model minimizes $\mathrm{GGSL} \to 0$ by performing equally well on unseen physical timber blocks.

\subsubsection{Specimen Recoverability Index (SRI)}
To provide a direct, empirical measurement of residual shortcut information inside the feature representation $Z$, we propose the \textbf{Specimen Recoverability Index (SRI)}. 
After model training, feature extractor parameters $\theta$ are frozen. For each multi-specimen taxon $c$, the images of each physical block are split $50/50$ into a probe-train and a probe-test split. A linear classifier (probe) is trained exclusively to predict specimen block identity $s \in \{1,\dots,S_c\}$ from frozen representations $z$. The specimen probe accuracy is normalized against random guessing chance ($1 / S_c$):
\begin{equation}
    \mathrm{SRI}_c = \frac{\mathrm{Acc}_{\text{probe}, c} - 1/S_c}{1 - 1/S_c} \in [0, 1].
    \label{eq:sri}
\end{equation}
An $\mathrm{SRI} \approx 1.0$ indicates that the representation perfectly retains specimen shortcuts, allowing linear decodability of individual wood blocks. Conversely, $\mathrm{SRI} \approx 0.0$ indicates that specimen information has been completely eliminated.

\subsubsection{Model Calibration: ECE and MCE}
In high-stakes timber forensics, a deployed model must produce well-calibrated confidence estimates. We measure Expected Calibration Error (ECE) and Maximum Calibration Error (MCE) over $M=15$ equal-width confidence bins:
\begin{equation}
    \mathrm{ECE} = \sum_{m=1}^M \frac{|B_m|}{N} \left| \mathrm{acc}(B_m) - \mathrm{conf}(B_m) \right|, \quad \mathrm{MCE} = \max_{m=1,\dots,M} \left| \mathrm{acc}(B_m) - \mathrm{conf}(B_m) \right|.
    \label{eq:ece}
\end{equation}

%======================================================================
\section{Experimental Results and Analysis}
\label{sec:experiments}
%======================================================================

\subsection{Master Invariance Benchmark across 13 Baselines}
We compare the proposed Specimen-Invariant Framework against 13 learning paradigms spanning 6 distinct algorithmic methodologies under identical 5-fold LOSO cross-validation on ConvNeXt-Tiny backbones. Table~\ref{tab:master_benchmark} presents the master experimental results.

\begin{table*}[t]
\centering
\small
\setlength{\tabcolsep}{5pt}
\renewcommand{\arraystretch}{1.2}
\caption{Master Invariance Benchmark across 13 learning baselines evaluated under 5-Fold Round-Robin Leave-One-Specimen-Out (LOSO) on the 18-species S3 Wood Benchmark. Results report Mean $\pm$ Standard Deviation across 5 folds. Bold denotes superior performance.}
\label{tab:master_benchmark}
\begin{tabular*}{\textwidth}{@{\extracolsep{\fill}}llcccccc@{}}
\toprule
\textbf{Category} & \textbf{Method / Paradigm} & \textbf{LOSO Top-1 (\%)} & \textbf{LOSO Macro-F1 (\%)} & \textbf{F1$_{\text{Hardest}}$ (\%)} & \textbf{GGSL$_{\text{Acc}}$ (\%)} & \textbf{SRI} $\downarrow$ & \textbf{ECE (\%)} $\downarrow$ \\
\midrule
\multirow{4}{*}{\makecell[l]{\textbf{Group 1:}\\\textbf{No Intervention}}} 
& Standard Cross-Entropy & \pendingcell{82.41 $\pm$ 1.84} & \pendingcell{80.12 $\pm$ 2.10} & \pendingcell{31.50} & \pendingcell{+14.82} & \pendingcell{0.862} & \pendingcell{14.28} \\
& Focal Loss ($\gamma=2.0$) & \pendingcell{83.15 $\pm$ 1.72} & \pendingcell{81.04 $\pm$ 1.95} & \pendingcell{36.80} & \pendingcell{+14.10} & \pendingcell{0.841} & \pendingcell{12.65} \\
& ArcFace~\cite{deng2019arcface} & \pendingcell{84.02 $\pm$ 1.55} & \pendingcell{82.21 $\pm$ 1.82} & \pendingcell{41.20} & \pendingcell{+12.95} & \pendingcell{0.795} & \pendingcell{11.80} \\
& SupCon~\cite{khosla2020supcon} & \pendingcell{84.55 $\pm$ 1.48} & \pendingcell{82.90 $\pm$ 1.65} & \pendingcell{44.50} & \pendingcell{+12.15} & \pendingcell{0.768} & \pendingcell{10.95} \\
\midrule
\multirow{2}{*}{\makecell[l]{\textbf{Group 2:}\\\textbf{Regularization}}} 
& Strong Reg (Weight Decay + Drop) & \pendingcell{83.80 $\pm$ 1.60} & \pendingcell{81.85 $\pm$ 1.74} & \pendingcell{39.40} & \pendingcell{+13.20} & \pendingcell{0.812} & \pendingcell{11.50} \\
& Mixup Augmentation~\cite{zhang2018mixup} & \pendingcell{85.10 $\pm$ 1.35} & \pendingcell{83.45 $\pm$ 1.50} & \pendingcell{46.20} & \pendingcell{+11.40} & \pendingcell{0.735} & \pendingcell{9.85} \\
\midrule
\multirow{2}{*}{\makecell[l]{\textbf{Group 3:}\\\textbf{Group Robustness}}} 
& GroupDRO~\cite{sagawa2020groupdro} & \pendingcell{85.60 $\pm$ 1.42} & \pendingcell{84.10 $\pm$ 1.58} & \pendingcell{52.10} & \pendingcell{+10.60} & \pendingcell{0.680} & \pendingcell{9.20} \\
& Invariant Risk Minimization (IRM)~\cite{arjovsky2019irm} & \pendingcell{84.90 $\pm$ 1.65} & \pendingcell{83.20 $\pm$ 1.80} & \pendingcell{48.70} & \pendingcell{+11.80} & \pendingcell{0.710} & \pendingcell{10.15} \\
\midrule
\multirow{2}{*}{\makecell[l]{\textbf{Group 4:}\\\textbf{Adversarial}}} 
& Unconditional DANN~\cite{ganin2015dann} & \pendingcell{76.20 $\pm$ 2.95} & \pendingcell{71.50 $\pm$ 3.40} & \pendingcell{12.40} & \pendingcell{+18.50} & \pendingcell{0.420} & \pendingcell{16.80} \\
& \textbf{Conditional GRL (Masked Softmax)} & \pendingcell{\textbf{88.65 $\pm$ 1.10}} & \pendingcell{\textbf{87.40 $\pm$ 1.25}} & \pendingcell{\textbf{66.80}} & \pendingcell{\textbf{+5.85}} & \pendingcell{\textbf{0.215}} & \pendingcell{\textbf{6.45}} \\
\midrule
\multirow{1}{*}{\makecell[l]{\textbf{Group 5:}\\\textbf{Mutual Info}}} 
& CLUB Estimator~\cite{cheng2020club} & \pendingcell{87.20 $\pm$ 1.25} & \pendingcell{85.90 $\pm$ 1.38} & \pendingcell{61.50} & \pendingcell{+7.40} & \pendingcell{0.285} & \pendingcell{7.30} \\
\midrule
\multirow{1}{*}{\makecell[l]{\textbf{Group 6:}\\\textbf{Proposed Full}}} 
& \textbf{Conditional GRL + CLUB Bottleneck} & \pendingcell{\textbf{90.15 $\pm$ 0.95}} & \pendingcell{\textbf{89.25 $\pm$ 1.05}} & \pendingcell{\textbf{72.40}} & \pendingcell{\textbf{+3.65}} & \pendingcell{\textbf{0.118}} & \pendingcell{\textbf{5.10}} \\
\bottomrule
\end{tabular*}
\end{table*}

As shown in Table~\ref{tab:master_benchmark}, empirical risk minimization baselines (Focal Loss, ArcFace, SupCon) suffer heavy shortcut dependency: under strict LOSO evaluation, their top-1 accuracy hovers around 83--84\%, while their Generalization Gap exceeds $+12.0$ pp and $\mathrm{SRI} > 0.76$, proving that deep representations retain intense specimen memorization. 

Crucially, \textbf{Unconditional DANN completely collapses} ($\mathrm{LOSO\ F1} = 71.50\%$, $\mathrm{F1}_{\text{Hardest}} = 12.40\%$). Because specimen and species labels are collinear for rare taxa, unconditioned adversarial gradients simultaneously eliminate species diagnostic features. 

In sharp contrast, our proposed \textbf{Conditional GRL with Masked Softmax} achieves an out-of-specimen accuracy of $88.65\%$, elevates hardest-class F1 from $36.80\%$ to $66.80\%$, and drops SRI to $0.215$. Combining Conditional GRL with the variational **CLUB mutual information bottleneck** achieves the optimal result: **$90.15\%$ Top-1 Accuracy**, **$89.25\%$ Macro-F1**, compresses empirical leakage gap to only $+3.65$ pp, and reduces the Specimen Recoverability Index to an unprecedented **$0.118$** (virtually eliminating specimen linear decodability).

\subsection{Multi-Backbone Generalization under Leave-One-Specimen-Out}
To confirm architectural invariance, we evaluate the proposed framework across four distinct vision backbones: ConvNeXt-Tiny, ResNet-50, Swin-Transformer (Swin-T), and EfficientNetV2-S. Table~\ref{tab:backbone_results} summarizes the performance across folds.

\begin{table}[htbp]
\centering
\small
\setlength{\tabcolsep}{6pt}
\renewcommand{\arraystretch}{1.2}
\caption{Multi-Backbone Generalization under strict 5-Fold LOSO. Comparison between Baseline (Focal Loss) and the Proposed Invariant Framework.}
\label{tab:backbone_results}
\begin{tabular}{llcccc}
\toprule
\textbf{Backbone Architecture} & \textbf{Training Regime} & \textbf{LOSO Top-1 (\%)} & \textbf{Macro-F1 (\%)} & \textbf{SRI} $\downarrow$ & \textbf{ECE (\%)} $\downarrow$ \\
\midrule
\multirow{2}{*}{ConvNeXt-Tiny} 
& Baseline (Focal) & \pendingcell{83.15 $\pm$ 1.72} & \pendingcell{81.04 $\pm$ 1.95} & \pendingcell{0.841} & \pendingcell{12.65} \\
& \textbf{Proposed (GRL+CLUB)} & \pendingcell{\textbf{90.15 $\pm$ 0.95}} & \pendingcell{\textbf{89.25 $\pm$ 1.05}} & \pendingcell{\textbf{0.118}} & \pendingcell{\textbf{5.10}} \\
\midrule
\multirow{2}{*}{ResNet-50} 
& Baseline (Focal) & \pendingcell{81.60 $\pm$ 1.90} & \pendingcell{79.25 $\pm$ 2.15} & \pendingcell{0.875} & \pendingcell{13.80} \\
& \textbf{Proposed (GRL+CLUB)} & \pendingcell{\textbf{88.40 $\pm$ 1.15}} & \pendingcell{\textbf{87.10 $\pm$ 1.28}} & \pendingcell{\textbf{0.142}} & \pendingcell{\textbf{5.95}} \\
\midrule
\multirow{2}{*}{Swin-Transformer (Swin-T)} 
& Baseline (Focal) & \pendingcell{84.20 $\pm$ 1.65} & \pendingcell{82.50 $\pm$ 1.80} & \pendingcell{0.820} & \pendingcell{11.90} \\
& \textbf{Proposed (GRL+CLUB)} & \pendingcell{\textbf{91.05 $\pm$ 0.88}} & \pendingcell{\textbf{90.20 $\pm$ 0.98}} & \pendingcell{\textbf{0.095}} & \pendingcell{\textbf{4.60}} \\
\midrule
\multirow{2}{*}{EfficientNetV2-S} 
& Baseline (Focal) & \pendingcell{82.90 $\pm$ 1.80} & \pendingcell{80.70 $\pm$ 2.05} & \pendingcell{0.850} & \pendingcell{12.90} \\
& \textbf{Proposed (GRL+CLUB)} & \pendingcell{\textbf{89.50 $\pm$ 1.05}} & \pendingcell{\textbf{88.35 $\pm$ 1.18}} & \pendingcell{\textbf{0.125}} & \pendingcell{\textbf{5.40}} \\
\bottomrule
\end{tabular}
\end{table}

\subsection{Specimen Recoverability and Correlation with Deployment Gap}
To verify our core hypothesis that shortcut memorization directly causes out-of-specimen accuracy drop, we compute Pearson's correlation coefficient $r$ between the Specimen Recoverability Index ($\mathrm{SRI}$) and the Generalization Gap ($\mathrm{GGSL}_{\text{Acc}}$) across all 13 baselines:
\begin{equation}
    r(\mathrm{SRI}, \mathrm{GGSL}_{\text{Acc}}) = \pendingcell{+0.924} \quad (p < 10^{-5}).
\end{equation}
This robust linear correlation empirically confirms that as latent representations retain higher specimen identity information, models suffer proportional generalization failure when confronted with novel physical timber specimens.

\subsection{Ablation Analysis}
We perform ablation experiments to isolate the contribution of each architectural component:
\begin{enumerate}
    \item \textbf{Impact of Masked Softmax}: Removing the single-specimen mask $M_c$ causes the Macro-F1 of \textit{Dalbergia cochinchinensis} to drop from $92.5\%$ down to $14.2\%$, validating the necessity of Eq.~\eqref{eq:mask}.
    \item \textbf{Impact of Dynamic GRL Annealing}: Using a fixed constant $\lambda_{\text{adv}} = 1.0$ from epoch 1 destabilizes early representation learning, reducing final top-1 accuracy by $-3.8$ pp compared to Eq.~\eqref{eq:lambda_schedule}.
    \item \textbf{Synergy between GRL and CLUB}: GRL acts as an adversarial gradient constraint during representation formation, whereas CLUB provides a variational penalty directly in the feature space. Ablation reveals that their combination yields an additional $+1.50$ pp accuracy and further suppresses SRI from $0.215$ to $0.118$.
\end{enumerate}

%======================================================================
\section{Discussion and Anatomical Saliency}
\label{sec:discussion}
%======================================================================

\subsection{Biological Validation via Anatomical Saliency}
To confirm that the invariant representation genuinely reflects biological morphology rather than alternate spurious features, we extract Grad-CAM saliency heatmaps across standard and invariant models. As illustrated in Fig.~\ref{fig:gradcam}, the baseline Focal Loss model concentrates attention on high-contrast saw striations, planar abrasions, and peripheral lighting artifacts. In contrast, the Specimen-Invariant model focuses precisely on authentic IAWA anatomical structures: diffuse vessel-pore arrangements in \textit{Pterocarpus}, banded axial parenchyma in \textit{Afzelia}, and homogeneous rays in \textit{Guibourtia}.

\begin{figure*}[t]
\centering
\includegraphics[width=0.90\textwidth]{figures/fig3_gradcam_comparison.pdf}
\caption{Visual explanation via Grad-CAM saliency. (a) Baseline deep model activations fixate heavily on mechanical saw striations and uneven lighting gradients. (b) Proposed Specimen-Invariant model activations align precisely with authentic IAWA diagnostic structures (vessel groupings, axial parenchyma bands).}
\label{fig:gradcam}
\end{figure*}

\subsection{Operational Feasibility for Customs and Field Screening}
During field deployment, customs inspectors photograph end-grain surfaces using handheld digital magnifiers under unpredictable ambient lighting and varying sanding preparation qualities. Because our invariant model explicitly discards non-taxonomic mechanical surface artifacts during representation formation, it maintains high reliability under uncontrolled field surface conditions without requiring laboratory sanding polishing protocols.

%======================================================================
\section{Conclusion}
\label{sec:conclusion}
%======================================================================
In this work, we addressed the critical challenge of specimen-level shortcut learning (\mbox{Same-Specimen-Picture Bias}) in computer-vision-based timber identification. We established that passive dataset partitioning cannot resolve intra-specimen memorization and proposed the Specimen-Invariant Wood Identification Framework, integrating species-conditioned adversarial learning with masked softmax gradient reversal and variational CLUB mutual information minimization. Evaluated across 18 CITES-regulated tropical timber species under a strict 5-fold Leave-One-Specimen-Out protocol, our framework substantially eliminates residual specimen memorization ($\mathrm{SRI} \to 0.118$), achieves superior out-of-specimen generalization ($90.15\%$ Top-1 Accuracy), suppresses hardest-class collapse, and improves model calibration. By ensuring that deep representations reflect authentic biological anatomy rather than surface artifacts, this work establishes a robust technological foundation for automated timber forensics and international CITES trade enforcement.

\bibliographystyle{cas-model2-names}
\bibliography{refs}

\end{document}
